\begin{table}[htbp]
\centering
\caption{Historical Spanish Exposure and Migration-Related Outcomes}
\begin{tabular}{l*{4}{c}}
\hline
            &\multicolumn{1}{c}{(1)}&\multicolumn{1}{c}{(2)}&\multicolumn{1}{c}{(3)}&\multicolumn{1}{c}{(4)}\\
            &\multicolumn{1}{c}{(1)}&\multicolumn{1}{c}{(2)}&\multicolumn{1}{c}{(3)}&\multicolumn{1}{c}{(4)}\\
\hline
Spanish share 1936-1955&       0.036\sym{***}&                     &       0.020\sym{**} &                     \\
            &     (0.009)         &                     &     (0.009)         &                     \\
[1em]
Spanish share 1956-1978&                     &       0.146\sym{***}&                     &       0.119\sym{***}\\
            &                     &     (0.025)         &                     &     (0.034)         \\
\hline
Observations&       1,515         &       1,515         &       1,463         &       1,463         \\
\hline
\end{tabular}
\captionsetup{justification = justified}
\caption*{\footnotesize \textit{Notes:} Each column reports the average marginal effect of the Spanish share in 1936–1955 from a separate logit regression for the outcome indicated in the column header. The dependent variable in columns (1) and (2) is an indicator for having a family member living abroad. The dependent variable in columns (3) and (4) is an indicator for intending to migrate to Spain, coded as 0 for individuals with no intention to migrate or with an intention to migrate to another country. Individual controls include age and male. Standard errors clustered at the municipality level in parentheses. * p<0.10, ** p<0.05, *** p<0.01. \textit{Sources}: LAPOP AmericasBarometer, 2023 round (dependent variables and individual controls); IPUMS Argentine Census 1970 (Spanish immigration share).}
\end{table}
